\scanpage{25}

% L1–4
Il y a enfin un sous topos $V$ \deleted{$\unsure{\Delta_1}$} qui n'est ni ouvert, ni fermé, \deleted{mais loc.} dont les pts sont $\Delta_0,\Delta_1$, et qui \deleted{\unsure{est}} est donné par
\[
 \begin{aligned}
 T_V(E)&=\text{cat. des sous objets de }\deleted{\unsure{X}}e_E\\
 &\text{(i.e. des objets }X\in\operatorname{Ob}E\text{ tels que }x\to x_e\text{ mono)}
 \end{aligned}
\]
% L5–7
et le sous-topos $\unsure{R^*|U^*|\mathrm{Ens}^{+}}\cap V$ \deleted{$\simeq\unsure{\varphi}$} topos ponctuel, point $\Delta_1$, tel que
\[
 \underset{\substack{\text{objets}\\\text{finaux}}}{T_{U\cap V}(E)}\subset E\qquad
 \text{i.e. }x\xrightarrow{\sim}e_E
\]

% L8
En résumé, on a $6$ sous-topos, avec diagramme d'implications
% L9–14, left diagram
\diagram{\begin{tikzcd}[ampersand replacement=\&,row sep=small,column sep=small]
 {} \& R\arrow[dl,dash]\arrow[dr,dash]\&{} \&{}\\
 R^*\arrow[dr,dash]\&{}\& V=\overline{(\varphi_0,\varphi_1)}\arrow[dl,dash]\arrow[dr,dash]\&{}\\
 {}\&R^*\cap V\arrow[dr,dash]\arrow[d,equal]\&{}\&\overline{(\varphi_0)}\arrow[dl,dash]\\
 {}\&\varphi_1\&\varnothing\&{}
\end{tikzcd}}
% L9–14, cancelled central diagram
\diagram{\begin{tikzcd}[ampersand replacement=\&,row sep=small,column sep=small,execute at end picture={\draw (\tikzcdmatrixname-1-2.north east)--(\tikzcdmatrixname-2-1.south west) (\tikzcdmatrixname-2-3.north east)--(\tikzcdmatrixname-3-2.south west) (\tikzcdmatrixname-2-3.north east)--(\tikzcdmatrixname-4-2.south west) (\tikzcdmatrixname-3-4.north east)--(\tikzcdmatrixname-4-3.south west);}]
 {}\&\deleted{R}\arrow[dl]\arrow[dr]\&{}\&{}\\
 \deleted{R^*}\arrow[dr]\&{}\&\deleted{\unsure{\Delta}=(\varphi_0,\varphi_1)}\arrow[dl]\arrow[dr]\&{}\\
 {}\&\deleted{R^*\cap\unsure{\Delta}=\varphi_1}\arrow[dr]\&{}\&\deleted{\varphi_0}\arrow[dl]\\
 {}\&{}\&\deleted{\varnothing}\&{}
\end{tikzcd}}
% L9–16, right diagram
\diagram{\begin{tikzcd}[ampersand replacement=\&,row sep=small,column sep=small]
 {}\&E\arrow[dl]\arrow[dr]\&{}\&{}\\
 \substack{\text{Objets}\\\text{couvrant }e_E}\arrow[dr]\&{}\&\substack{\text{sous-objets}\\\text{de }e_E}\arrow[dl]\arrow[dr]\&{}\\
 {}\&\substack{\text{objets}\\\text{finaux}}\arrow[dr]\&{}\&\text{objets vides}\arrow[dl,dash]\\
 {}\&{}\&\substack{\varnothing\ \text{sauf si }E=\mathrm{Top}(\varnothing)\\\text{où toutes les}\\\text{cat envisagées}\\\text{sont }=E=\\\text{cat ponctuelle}}\&{}
\end{tikzcd}}

% L13–17, left
\uline{NB} \deleted{\unsure{RR}} $V$ et $R^*\cap V$ n'ont pas de complémentaire mais ils ont les complémentaires faibles $\varnothing$ et $\varphi_0$

% L18–19, left
$V$\inserted{$\unsure{?}$} représente le \inserted{$?$} foncteur
\[
 E\longmapsto\unsure{\mathfrak P|\mathscr P|\mathcal P}(E)=\mathfrak P(e_E)
\]

% L16–21, right diagram
\diagram{\begin{tikzcd}[ampersand replacement=\&,row sep=small,column sep=small]
 {}\&\mathrm{Ens}\arrow[dl,dash]\arrow[dr,dash]\&{}\&{}\\
 \text{Ens. non vides}\arrow[dr,dash]\&{}\&\substack{\text{ens. de}\\\text{card }\leq 1}\arrow[dl,dash]\arrow[dr,dash]\&{}\\
 {}\&\substack{\text{ens.}\\\text{ponctuels}}\arrow[dr,dash]\&{}\&\{\varnothing\}\arrow[dl,dash]\\
 {}\&{}\&\varnothing\&{}
\end{tikzcd}}

% L21–23
\uline{NB}
\[
 \left\{\begin{array}{l}
 \text{Pas d'objets loc. constants, ou constants}\\
 \text{Pas d'objets \deleted{constants} qui \unsure{après un pr. sup|après une pr. sup} ont}\\
 \hspace{3em}\text{\deleted{\unsure{qui après}} (resp. id sur }e\text{)}\\
 \text{Pas d'objets constructibles}
 \end{array}\right.
\]

% L24–29, left (sideways)
% margin note written sideways
\margin{\unsure{C'est} donc des structures alg. définissables par $\underleftarrow{\lim}$ finies et même par produits finis, et sauf les deux derniers \unsure{$\Longleftarrow$} par \unsure{l'un des deux|l'un des trois} ---}

% L24–29
\uline{Question} Topos classifiant et sous-topos pour la structure
\[
 \begin{aligned}
 &\underbrace{X,Y}\quad\underbrace{f:X\to Y\ \text{sans plus}},\ \text{\unsure{ou|où}}\quad
 \underbrace{f:x-\unsure{x|X}\ \text{sans plus}}_{B_{\mathbb N^{\times}}},\quad \underbrace{f:e\to X}\\
 &\deleted{\unsure{\varphi}}\qquad X\times Y\to Z,\quad X\times Y\to X\\
 &\hspace{5em}X\times X\longrightarrow X\qquad !\\
 &\hspace{5em}X\times X\rightrightarrows x\qquad !!
 \end{aligned}
\]
% archive sheet number: 24
